\chapter{Stratégies d'investissement quantitatives}\label{ch:s2:quantitative-investment-strategies}

Une banque peut regrouper une stratégie de carry ou de momentum sous forme d'indice, la licencier et vendre des swaps dessus. Le backtest qui sert à vendre l'indice est rédigé avant la
mise en conditions réelles de l'indice, et ses rendements après lancement sont plus faibles. Suhonen, Lennkh et Perez ont étudié 215 stratégies de ce type, mises au point et promues par de grandes banques
d'investissement mondiales et ont constaté une baisse médiane de 73~\% des ratios de Sharpe du backtest à la période en conditions réelles. Dans les banques synthétiques de ce chapitre, chacune lançant
les trois meilleurs sur vingt indices candidats, les indices lancés affichent un ratio de Sharpe de 0.51 au backtest et de 0.16 en conditions réelles après frais, et plus d'un tiers
d'entre eux perdent de l'argent en conditions réelles. La construction est assurée par \texttt{firm.qis}.

\section{Indices fondés sur des règles comme produits}

\begin{definition}[Stratégie d'investissement quantitative]\label{def:s2:quantitative-investment-strategies:qis}
Une \emph{stratégie d'investissement quantitative}\index{stratégie d'investissement quantitative} (QIS) est une stratégie systématique fondée sur des règles, publiée par une banque sous forme
d'indice (un indice de stratégie systématique, livre~5, chapitre~19), dotée d'un règlement de marché, d'un agent de calcul et d'un backtest, et à laquelle les clients accèdent par des swaps,
des notes ou des fonds plutôt qu'en la faisant tourner eux-mêmes.
\end{definition}

Un indice QIS est une stratégie entourée d'un contrat~: les règles fixent ce que l'indice détient, quand il se rééquilibre et quels frais il prélève, de sorte que le client
peut acheter une exposition au carry, au momentum, à la value ou à une prime de volatilité sans équipe pour la faire tourner. La banque perçoit des frais et la marge sur les swaps, et elle
décide quels indices lancer. Cette décision se prend sur la base de backtests.

\texttt{firm.qis} fournit à 2\,000 équipes produit vingt indices candidats chacune (le \cref{lst:s2:quantitative-investment-strategies:launch}). Les ratios de Sharpe réels des
candidats sont tirés autour de 0.10, avec un écart-type de 0.15, et leurs erreurs d'estimation sont corrélées à 0.5, puisqu'il s'agit de variantes
de quelques idées. Chaque équipe fait un backtest sur dix ans, lance ses trois meilleurs indices et les conserve pendant cinq ans, avec une cible de volatilité de 10~\%. Les
indices en conditions réelles supportent 0.5~\% par an de frais d'indice et de swap, et perdent 0.2~\% par an du fait du rééquilibrage prévisible. Les ratios de Sharpe réels sont calibrés pour que la
médiane d'érosion corresponde aux 73~\% que rapportent Suhonen et ses coauteurs~; tout le reste découle de la sélection.

\section{Enveloppes de stratégie et frais}

\begin{definition}[Enveloppe de stratégie]\label{def:s2:quantitative-investment-strategies:wrapper}
Une \emph{enveloppe de stratégie}\index{enveloppe de stratégie} (strategy wrapper) est la forme juridique et économique sous laquelle un client détient un indice de stratégie~: un total return swap,
une note structurée, un certificat ou un fonds, chacun avec ses propres frais, sa contrepartie, son traitement prudentiel et sa liquidité.
\end{definition}

Avec une cible de volatilité de 10~\%, 0.7~\% par an de frais et de fuites coûtent 0.07 de ratio de Sharpe. C'est une part importante d'un ratio de Sharpe réel de 0.1 à 0.3,
et il est invisible dans un backtest réalisé avant que l'indice n'existe~: le backtest ne comporte aucun frais, parce que personne n'en payait. L'enveloppe ajoute ses propres conditions~:
un swap comporte un risque de contrepartie et un spread de financement, et une note ajoute le spread de crédit de l'émetteur (livre~5, chapitre~19).

\section{Couverture de l'indice}

Une banque qui verse à un client la performance d'un indice au moyen d'un swap se couvre en détenant les positions de l'indice. Le règlement de marché indique à tous quand l'indice
se rééquilibrera et à peu près ce qu'il négociera, de sorte que d'autres peuvent passer avant lui, et que l'indice achète un peu plus haut et vend un peu plus bas qu'il ne le ferait
sans être observé. Le chapitre impute cela comme une fuite de 0.2~\% par an, supportée seulement en conditions réelles, puisqu'un backtest négocie au cours de clôture historique. La même
prévisibilité permet à la banque de compenser les transactions de l'indice avec ses autres flux (chapitre~25), ce qui réduit son propre coût de couverture mais pas celui du client.

\section{Conditions réelles contre backtest}

\begin{definition}[Érosion après backtest de l'indice]\label{def:s2:quantitative-investment-strategies:decay}
L'\emph{érosion après backtest de l'indice}\index{erosion@érosion après backtest de l'indice} (index backtest decay) est la baisse du ratio de Sharpe entre le backtest d'un indice de stratégie réalisé avant son lancement et celui de ses
rendements en conditions réelles, nets de frais, mesurés sur l'ensemble des indices lancés.
\end{definition}

\begin{center}
\small
\setlength{\tabcolsep}{4pt}%
\begin{tabular}{@{}lcc@{}}
\toprule
Ratios de Sharpe, moyenne sur les indices lancés & meilleurs trois sur vingt & meilleur sur vingt\\
\midrule
backtest (dix ans) & 0.51 & 0.62\\
réel & 0.22 & 0.26\\
en conditions réelles, avant frais (cinq ans) & 0.23 & 0.27\\
en conditions réelles, après frais et fuites & 0.16 & 0.20\\
érosion médiane, du backtest aux conditions réelles après frais & 72.2~\% & 69.4~\%\\
\bottomrule
\end{tabular}
\end{center}

Le ratio de Sharpe en conditions réelles est proche du ratio réel~: les rendements en conditions réelles sont honnêtes, simplement bruités. Le backtest ne l'est pas, parce qu'il a été choisi pour être élevé. La sélection
a bien trouvé de meilleures stratégies, avec un ratio de Sharpe réel de 0.22 contre 0.10 pour le candidat moyen, mais elle leur a aussi attribué du bruit. Après
frais, 36.6~\% des indices lancés affichaient un ratio de Sharpe négatif en conditions réelles. Le meilleur indice d'une équipe montre à quoi cela ressemble de l'intérieur
(la \cref{fig:s2:quantitative-investment-strategies:path})~: un ratio de Sharpe de 0.71 au backtest et un gain de 71~\% sur dix ans, puis cinq années en conditions réelles avec un
ratio de Sharpe de $-0.07$ avant frais et une perte de 7.1~\% après frais.

\begin{omfigure}
\begin{tikzpicture}
\begin{axis}[width=11cm,height=5.2cm,xlabel={\small années depuis le lancement},ylabel={\small rendement cumulé (\%)},tick label style={font=\footnotesize},
  xmin=-10,xmax=5,ymin=-20,ymax=90,ytick={-20,0,20,40,60,80},grid=major,grid style={gray!25},table/col sep=comma]
\fill[gray!15] (axis cs:0,-20) rectangle (axis cs:5,90);
\addplot[omThm,thick] table[x=year,y=cum]{figdata/strategies-2/27-quantitative-investment-strategies/path.csv};
\node[font=\scriptsize,anchor=north west] at (axis cs:-9.6,85) {backtest};
\node[font=\scriptsize,anchor=north west] at (axis cs:0.3,85) {réel};
\end{axis}
\end{tikzpicture}

\omcaption{Le meilleur indice d'une équipe produit synthétique~: somme cumulée des rendements quotidiens à une cible de volatilité de 10~\%, sur son backtest de dix ans et ses
cinq premières années en conditions réelles, nets de frais et de fuites (zone grisée). Données~: \texttt{s2\_qis.path}.}\label{fig:s2:quantitative-investment-strategies:path}
\end{omfigure}

Le ratio de Sharpe déflaté (livre~4, chapitre~12) demande si le meilleur de $N$ backtests bat ce que montrerait le meilleur de $N$ backtests sans valeur. Bailey et
Lopez de Prado l'ont conçu pour corriger précisément cette sélection. Appliqué au meilleur indice de chaque équipe, avec vingt essais et la dispersion de leurs propres
ratios de Sharpe de backtest, il en valide 5.9~\% au niveau de confiance de 95~\%. Ceux qui passent ont affiché 1.13 au backtest et 0.32 en conditions réelles avant frais~; ceux qui échouent,
à 0.59 et 0.27. Le test a raison de se montrer sceptique, mais dix ans de backtest ne permettent pas de distinguer un ratio de Sharpe de 0.2 d'un ratio de 0.3~: la discipline qu'il apporte
tient dans les backtests qu'il empêche d'être vendus comme excellents, et non dans le choix des gagnants.

Suhonen et ses coauteurs ont aussi constaté que les stratégies les plus complexes ont perdu plus de 30~points de pourcentage supplémentaires par rapport aux plus simples. La complexité signifie plus de
candidats testés, ouvertement ou non. Dans le modèle, choisir le meilleur parmi davantage de candidats gonfle le backtest bien plus que le résultat en conditions réelles
(la \cref{fig:s2:quantitative-investment-strategies:complexity})~: de 2 à 100 candidats, le meilleur backtest passe de 0.26 à 0.76 et le ratio de Sharpe
en conditions réelles après frais passe de 0.07 à 0.24, si bien que l'écart entre les deux passe de 0.18 à 0.52. McLean et Pontiff ont relevé une érosion comparable chez les prédicteurs d'actions
publiés, dont les rendements étaient inférieurs de 26~\% hors échantillon (out-of-sample) et de 58~\% après publication.

\begin{omfigure}
\begin{tikzpicture}
\begin{axis}[width=11cm,height=5.2cm,xmode=log,xlabel={\small indices candidats testés (échelle logarithmique)},ylabel={\small ratio de Sharpe du meilleur},
  tick label style={font=\footnotesize},xmin=2,xmax=100,ymin=0,ymax=0.85,xtick={2,5,10,20,50,100},xticklabels={2,5,10,20,50,100},grid=major,
  grid style={gray!25},table/col sep=comma,legend style={font=\scriptsize,at={(0.03,0.97)},anchor=north west}]
\addplot[omThm,thick,mark=*,mark size=1.2pt] table[x=candidates,y=backtest]{figdata/strategies-2/27-quantitative-investment-strategies/complexity.csv};
\addplot[omDef,thick,mark=square*,mark size=1.2pt] table[x=candidates,y=net]{figdata/strategies-2/27-quantitative-investment-strategies/complexity.csv};
\legend{backtest,réel après frais}
\end{axis}
\end{tikzpicture}

\omcaption{Le meilleur des $N$ indices candidats~: ratio de Sharpe moyen du backtest et ratio de Sharpe moyen en production après frais et fuites, sur 2\,000 équipes
de produits synthétiques, en fonction du nombre de candidats testés. Données~: \texttt{s2\_qis.complexity}.}\label{fig:s2:quantitative-investment-strategies:complexity}
\end{omfigure}

\section{Fiches de stratégie}

\begin{strategyfile}{Swap d'indice de primes de risque}\label{strat:s2:quantitative-investment-strategies:rp}
\sfield{Qui paie, et pourquoi} Clients qui veulent des primes systématiques sans les gérer eux-mêmes~; la banque perçoit des commissions et la marge du swap.
\sfield{Instruments et plateformes} Total return swaps sur les indices de la banque~; les contrats à terme et actions sous-jacents à l'indice servent de couverture.
\sfield{Signal} Le règlement de l'indice.
\sfield{Dimensionnement et exécution} Couverture par la gestion de l'indice~; compenser ses transactions avec les autres flux.
\sfield{Coûts} Fuites liées au rééquilibrage~; financement~; capital.
\sfield{Mode de disparition} Rendements en production très inférieurs au backtest~; les clients partent.
\sfield{Horizon, capacité, infrastructure} Plusieurs années~; un agent de calcul de l'indice et un desk de couverture.
\sfield{Backtest en toute honnêteté} Compter chaque candidat testé~; déduire les frais et les fuites~; présenter le ratio de Sharpe déflaté.
\sfield{Sources} Suhonen, Lennkh et Perez (2017)~; ce chapitre~: backtest 0.51, en production après frais 0.16.
\end{strategyfile}

\begin{strategyfile}{Indice de carry de volatilité}\label{strat:s2:quantitative-investment-strategies:volcarry}
\sfield{Qui paie, et pourquoi} Acheteurs d'assurance (chapitre~1)~; l'indice la vend de façon systématique.
\sfield{Instruments et plateformes} Options sur indices ou swaps de variance, roulés selon une règle.
\sfield{Signal} Le règlement de marché~: échéance, strike et taille.
\sfield{Dimensionnement et exécution} Notionnel fixe ou ciblé en volatilité.
\sfield{Coûts} Spreads sur options~; rolls prévisibles négociés à l'avance.
\sfield{Mode de disparition} Une envolée de la volatilité~; des rolls surpeuplés.
\sfield{Horizon, capacité, infrastructure} Quelques mois~; accès au marché des options.
\sfield{Backtest en toute honnêteté} Inclure les krachs et les coûts de roll réellement payés.
\sfield{Sources} Suhonen, Lennkh et Perez (2017) ont trouvé que les expositions à la volatilité actions figuraient parmi les plus robustes en production.
\end{strategyfile}

\begin{strategyfile}{Panier multi-actifs de primes de risque}\label{strat:s2:quantitative-investment-strategies:basket}
\sfield{Qui paie, et pourquoi} Diversification entre les primes de carry, de momentum, de value et de volatilité.
\sfield{Instruments et plateformes} Plusieurs indices combinés dans un même swap ou fonds.
\sfield{Signal} Règles d'allocation (risque égal).
\sfield{Dimensionnement et exécution} Cible de volatilité au niveau du panier.
\sfield{Coûts} Frais et fuites de chaque indice.
\sfield{Mode de disparition} Des primes qui échouent ensemble~; corrélation cachée.
\sfield{Horizon, capacité, infrastructure} Plusieurs années.
\sfield{Backtest en toute honnêteté} Les résultats en production de chaque composante, et non son backtest.
\sfield{Sources} Aucun chiffre de performance vérifié.
\end{strategyfile}

\begin{strategyfile}{Indice défensif en superposition}\label{strat:s2:quantitative-investment-strategies:defensive}
\sfield{Qui paie, et pourquoi} Investisseurs qui paient pour réduire les drawdowns (chapitre~7).
\sfield{Instruments et plateformes} Indices de tendance ou fondés sur des puts, superposés aux positions actions.
\sfield{Signal} La règle de tendance ou de protection du règlement de marché.
\sfield{Dimensionnement et exécution} Dimensionné sur le portefeuille qu'il protège.
\sfield{Coûts} Le coût de carry de la protection~; frais.
\sfield{Mode de disparition} Une protection qui se vide de sa valeur pendant des années avant un krach qu'elle rate.
\sfield{Horizon, capacité, infrastructure} Plusieurs années.
\sfield{Backtest en toute honnêteté} Le juger comme une assurance~: coût rapporté aux pertes évitées, en production.
\sfield{Sources} Aucun chiffre de performance vérifié.
\end{strategyfile}

\section{Tutoriel~: rédigé avant le lancement}\label{tut:s2:quantitative-investment-strategies:tut}

\begin{tutorial}
\textbf{Objectif.} Simuler de nombreuses équipes de produits qui choisissent des indices à partir de backtests, appliquer frais et fuites, comparer la production au backtest, et tester le choix
avec le ratio de Sharpe déflaté. \textbf{État final~:} le tableau et les deux figures.
\begin{tutsteps}
\item \textbf{Équipes et lancements}.
\omcode{code/firm/qis/firm_qis.py}{48}{70}{Candidats aux erreurs d'estimation corrélées~; les k meilleurs sont lancés et soumis aux coûts.}\label{lst:s2:quantitative-investment-strategies:launch}
\item \textbf{Le verdict déflaté}.
\omcode{code/firm/qis/firm_qis.py}{73}{90}{Le meilleur de chaque équipe comparé au meilleur attendu de candidats sans valeur.}\label{lst:s2:quantitative-investment-strategies:dsr}
\item \textbf{Exécuter} \texttt{launched()}, \texttt{verdicts()}, \texttt{complexity()}, \texttt{path()} et \texttt{fig\_qis.py}.
\end{tutsteps}
\textbf{Pour aller plus loin.} Laisser les équipes cacher les candidats échoués (ne compter que les survivants) et constater que le test déflaté se laisse abuser~; ajouter des frais qui dépendent du
backtest~; laisser les clients partir après deux mauvaises années en production.
\end{tutorial}

\section{Construction~: stratégies d'investissement quantitatives}\label{bld:s2:quantitative-investment-strategies:qis}

\begin{build}
\bfield{Rôle} Sélection d'indices par backtest, frais et fuites, production face au backtest, et verdict du ratio de Sharpe déflaté.
\bfield{Interface} \texttt{QISConfig(\dots)}, \texttt{simulate\_teams(cfg)}, \texttt{launch(sim, cfg, k)}, \texttt{deflated(sim, cfg)},
\texttt{example\_path(cfg, team)}.
\bfield{Règles} Les ratios de Sharpe estimés sont la vérité plus une erreur d'échantillonnage de $1/\sqrt{T}$~; frais et fuites seulement en production~; le ratio de Sharpe
déflaté du livre~4.
\bfield{Tests d'acceptation} \texttt{code/firm/qis/tests/}~: l'ampleur de l'erreur d'échantillonnage~; la sélection gonfle le backtest, pas la vérité~; probabilités
comprises dans leur intervalle et le ratio de Sharpe du chemin d'exemple.
\bfield{Pour aller plus loin} Rendements non gaussiens~; candidats de complexité différente~; essais cachés.
\end{build}

\begin{omsources}
\item A.~Suhonen, M.~Lennkh et F.~Perez, «~Quantifying backtest overfitting in alternative beta strategies~», \emph{Journal of Portfolio Management}
43(2), 2017.
\item D.~H.~Bailey et M.~Lopez de Prado, «~The deflated Sharpe ratio: correcting for selection bias, backtest overfitting, and non-normality~»,
\emph{Journal of Portfolio Management} 40(5), 2014.
\item R.~D.~McLean et J.~Pontiff, «~Does academic research destroy stock return predictability~?~», \emph{Journal of Finance} 71(1), 2016.
\end{omsources}

\section{Exercices}

\begin{exercise}[$\star$]\label{exo:s2:quantitative-investment-strategies:1}
Un indice vise une volatilité de 10~\% et coûte 0.7~\% par an en frais et fuites. Calculer la part du ratio de Sharpe absorbée par ces coûts.
\end{exercise}

\begin{exercise}[$\star$]\label{exo:s2:quantitative-investment-strategies:2}
Un backtest sur dix ans affiche un ratio de Sharpe de 0.5. Calculer son erreur type sous l'hypothèse de rendements normaux.
\end{exercise}

\begin{exercise}[$\star$]\label{exo:s2:quantitative-investment-strategies:3}
Backtest 0.51, en production après frais 0.16~: calculer l'érosion.
\end{exercise}

\begin{exercise}[$\star\star$]\label{exo:s2:quantitative-investment-strategies:4}
Expliquer pourquoi le ratio de Sharpe en production est proche de la valeur vraie, alors que celui du backtest le dépasse largement.
\end{exercise}

\begin{exercise}[$\star\star$]\label{exo:s2:quantitative-investment-strategies:5}
Expliquer pourquoi les stratégies plus complexes s'érodent davantage.
\end{exercise}

\begin{exercise}[$\star\star$]\label{exo:s2:quantitative-investment-strategies:6}
Expliquer pourquoi un calendrier de rééquilibrage publié coûte à l'indice.
\end{exercise}

\begin{exercise}[$\star\star\star$]\label{exo:s2:quantitative-investment-strategies:7}
\emph{Programmation.} Exécuter \texttt{complexity()}. Comment le meilleur backtest et le résultat en production évoluent-ils lorsque le nombre de candidats passe de 2 à 100,
et pourquoi l'érosion médiane bouge-t-elle à peine~?
\end{exercise}

\begin{exercise}[$\star\star\star$]\label{exo:s2:quantitative-investment-strategies:8}
\emph{Trouver la faille.} «~L'indice a réussi le test de Sharpe déflaté, donc son ratio de Sharpe en production sera proche de celui du backtest.~»
\end{exercise}

\section{Problème~: écrit avant le lancement}

\begin{problem}[{Problème du week-end --- stratégies d'investissement quantitatives}]\label{pb:s2:quantitative-investment-strategies:1}
Les équipes produit synthétiques du chapitre et le dossier public.

\medskip\noindent\textbf{Partie I --- Le produit.}
\begin{enumerate}
\item Définir une \omterm{def:s2:quantitative-investment-strategies:qis}{stratégie d'investissement quantitative} (QIS) et une \omterm{def:s2:quantitative-investment-strategies:wrapper}{enveloppe de stratégie} (wrapper).
\item Comment la banque tire-t-elle un revenu d'un QIS~?
\item Décrire les équipes synthétiques et les candidats.
\item Qu'est-ce qui est calibré, et sur quoi~?
\end{enumerate}

\medskip\noindent\textbf{Partie II --- Coûts.}
\begin{enumerate}[resume]
\item Que coûtent les frais et les fuites en ratio de Sharpe~?
\item Pourquoi sont-ils absents du backtest~?
\item Comment la banque couvre-t-elle un swap sur indice~?
\item Que coûte la prévisibilité à l'indice~?
\end{enumerate}

\medskip\noindent\textbf{Partie III --- En production face au backtest.}
\begin{enumerate}[resume]
\item Définir l'\omterm{def:s2:quantitative-investment-strategies:decay}{érosion après backtest de l'indice}.
\item Donner le tableau des ratios de Sharpe du backtest, du ratio vrai et de la production.
\item Quelle part des indices lancés perd en production~?
\item Que dit le ratio de Sharpe déflaté~?
\end{enumerate}

\medskip\noindent\textbf{Partie IV --- Le verdict.}
\begin{enumerate}[resume]
\item Énoncer le \emph{résultat nommé}~: l'écart entre les ratios de Sharpe du backtest et de la production des indices sélectionnés.
\item Que trouvent Suhonen, Lennkh et Perez~?
\item Que fait la complexité dans le modèle~?
\item Que trouvent McLean et Pontiff~?
\item Comment un client doit-il lire le backtest d'un QIS~?
\item Quelle fiche de stratégie Suhonen et ses coauteurs jugent-ils le plus robuste~?
\item Comment ce chapitre se rapporte-t-il au chapitre du livre~7 sur le surapprentissage~?
\item En une phrase~: que vend un desk QIS~?
\end{enumerate}
\end{problem}

\section{Questions d'entretien}

\begin{interviewq}[$\star$ \iqroles{trader}]\label{iq:s2:quantitative-investment-strategies:1}
Qu'est-ce qu'un indice QIS, et comment une banque couvre-t-elle un swap sur cet indice~?
\end{interviewq}

\begin{interviewq}[$\star\star$ \iqroles{chercheur}]\label{iq:s2:quantitative-investment-strategies:2}
Une équipe produit vous montre un backtest de dix ans avec un ratio de Sharpe de 1. Que demandez-vous~?
\end{interviewq}

\begin{interviewq}[$\star\star$ \iqroles{trader}]\label{iq:s2:quantitative-investment-strategies:3}
D'autres traders anticipent le rééquilibrage mensuel de votre indice. Que pouvez-vous modifier dans le règlement de marché~?
\end{interviewq}

\begin{interviewq}[$\star\star$ \iqroles{risque}]\label{iq:s2:quantitative-investment-strategies:4}
Quels risques la banque conserve-t-elle lorsqu'elle vend des swaps sur ses propres indices~?
\end{interviewq}

\begin{interviewq}[$\star\star$ \iqroles{développeur}]\label{iq:s2:quantitative-investment-strategies:5}
Que doit garantir le système de l'agent de calcul de l'indice~?
\end{interviewq}

\begin{interviewq}[$\star\star\star$ \iqroles{chercheur}]\label{iq:s2:quantitative-investment-strategies:6}
Les ratios de Sharpe estimés des candidats sont la vérité plus un bruit indépendant d'écart type $s$, et les vérités sont égales. Montrez que le
backtest attendu du meilleur parmi $N$ dépasse son espérance en production d'environ $s\sqrt{2\ln N}$ pour $N$ grand.
\end{interviewq}
